\chapter{Документирование и тестирование}
\label{ch:10}

Вы умеете писать классы, собирать проект и поднимать веб-приложение. Код работает~--- когда запускаете его вы и вводите данные, которые сами же придумали. Проблемы начинаются позже: через месяц вы не помните, почему метод возвращает \texttt{null}, а не бросает исключение; через два~--- правка одной строки ломает совсем другой экран. Представьте, что вы купили сложную кофемашину: к ней прилагаются инструкция и заводской протокол испытаний. Без инструкции вы будете тыкать кнопки наугад; без испытаний узнаете о браке, когда кипяток окажется на полу. В программировании эти две вещи называются \term{документированием} и \term{тестированием}. Прочитав главу, вы научитесь собирать готовую справку из комментариев, писать тесты на JUnit~5, проверять исключения и наборы данных, изолировать класс от базы и сети с помощью Mockito и тестировать Spring-приложение на трёх уровнях.

\section{Комментарии и документирующий комментарий}

Три синтаксические формы комментария разобраны в главе~\ref{ch:01}: однострочный \texttt{//}, блочный \texttt{/* $\ldots$ */} и документирующий \texttt{/** $\ldots$ */}. Компилятор игнорирует все три~--- в байт-код они не попадают,~--- но третий умеет читать отдельная утилита \texttt{javadoc} из состава JDK, и потому он занимает особое место.

Разница не только в лишней звёздочке. Документирующий комментарий ставится непосредственно перед объявлением класса, интерфейса, метода, конструктора или поля: внутри тела метода он документацией не станет. Он делится на описание и блок тегов, а внутри разрешена HTML-разметка. Бытовое сравнение: однострочный комментарий~--- записка на холодильнике «молоко в дверце», она для того, кто уже стоит на этой кухне; Javadoc~--- инструкция к холодильнику, её читает человек, который к вашей кухне не подходил.

\begin{oshibka}
\textbf{Частая ошибка: комментарий, дублирующий код.} Строки вроде \texttt{i = i + 1; // увеличиваем i на единицу} или «Геттер для поля name» над \texttt{getName()} не просто бесполезны~--- они опасны. Код меняется, комментарий остаётся, через полгода они противоречат друг другу, а человек верит комментарию. Хороший комментарий отвечает не на вопрос «что делает код» (это видно из кода), а на вопросы «почему именно так» и «что нужно знать снаружи».
\end{oshibka}

\section{Javadoc: структура документирующего комментария}

\term{Javadoc}~--- стандартный генератор документации из исходного кода, входящий в состав JDK. Он читает файлы \texttt{.java}, извлекает документирующие комментарии вместе с сигнатурами классов и методов и формирует набор связанных HTML-страниц; именно так сделана официальная документация по стандартной библиотеке Java. Аналогия: Javadoc~--- типография; вы пишете рукопись на полях кода, а она собирает книжку с оглавлением и указателем.

Комментарий состоит из трёх частей, идущих строго в этом порядке. \term{Краткое описание} (\eng{summary})~--- всё до первой точки с пробелом после неё; именно этот текст попадает в сводные таблицы «Method Summary». Отсюда правило: не ставьте точку в середине первого предложения, иначе краткое описание оборвётся; если точка неизбежна (сокращение «т.\,д.»), задайте границу тегом \texttt{\{@summary $\ldots$\}}. \term{Подробное описание}~--- всё до первого блочного тега; здесь допустима разметка \texttt{<p>}, \texttt{<ul>}, \texttt{<pre>}. \term{Блок тегов}~--- строки, начинающиеся с символа \texttt{@}; как только он начался, обычный текст закончился.

Собрав всё вместе, получим класс из листинга~\ref{lst:10-1}~--- сквозной пример главы. Обратите внимание на три детали: поля тоже документируются коротким однострочным комментарием; пример кода обёрнут в \texttt{<pre>} и \texttt{\{@code $\ldots$\}}, чтобы угловые скобки не сломали HTML; тег \texttt{@throws} описывает не любое возможное исключение, а то, которое является частью контракта.

\begin{listing}[H]
\begin{minted}{java}
package ru.fa.library.service;

/**
 * Калькулятор штрафов за несвоевременный возврат книг.
 *
 * <p>Штраф начисляется за каждый день просрочки по ставке
 * {@link #DAILY_RATE} и ограничен значением {@link #MAX_FINE}.
 *
 * <p>Пример: <pre>{@code
 * double fine = new FineCalculator().calculate(10);  // 50.0
 * }</pre>
 *
 * @author А. Эбрахим
 * @version 1.2
 * @since 1.0
 */
public class FineCalculator {

    /** Штраф за один день просрочки, в рублях. */
    public static final double DAILY_RATE = 5.0;

    /** Максимальный штраф за одну книгу, в рублях. */
    public static final double MAX_FINE = 500.0;

    /**
     * Вычисляет штраф за указанное число дней просрочки.
     *
     * <p>Результат равен произведению дней на ставку, но не
     * превышает {@link #MAX_FINE}; для нуля вернёт {@code 0.0}.
     *
     * @param daysOverdue дней просрочки, неотрицательное число
     * @return сумма штрафа в рублях
     * @throws IllegalArgumentException если дней меньше нуля
     */
    public double calculate(int daysOverdue) {
        if (daysOverdue < 0) {
            throw new IllegalArgumentException(
                    "Число дней отрицательное: " + daysOverdue);
        }
        return Math.min(daysOverdue * DAILY_RATE, MAX_FINE);
    }
}
\end{minted}
\caption{Полностью документированный класс}
\label{lst:10-1}
\end{listing}

Правила хорошего описания коротки. Метод описывайте глаголом в третьем лице («Вычисляет штраф»), класс~--- существительным. Первое предложение должно быть самодостаточным: его читают в таблице методов, вне контекста. Документируйте контракт, а не реализацию: диапазон результата, поведение на границах, что будет при \texttt{null}, бросается ли исключение.

\section{Дескрипторы Javadoc}

\term{Дескриптор} (тег) Javadoc~--- метка с символа \texttt{@}, помечающая фрагмент документации как имеющий определённый смысл. Аналогия~--- почтовая накладная: на коробке можно написать что угодно, но адрес, отправитель и вес должны стоять в отведённых графах, иначе автоматика на сортировочном узле их не прочитает. \term{Блочные} дескрипторы пишутся с новой строки в блоке тегов (табл.~\ref{tab:10-1}), \term{строчные} (инлайновые)~--- в фигурных скобках прямо внутри текста.

\begin{table}[htbp]
\caption{Блочные дескрипторы Javadoc}
\label{tab:10-1}
\centering
\begin{tabular}{L{4.4cm}L{3.6cm}L{7cm}}
\toprule
Дескриптор & Где применяется & Назначение \\
\midrule
\texttt{@param} имя описание & метод, обобщённый тип, запись & параметр метода, параметр типа или компонент записи \\
\texttt{@return} описание & метод & возвращаемое значение \\
\texttt{@throws}, \texttt{@exception} & метод, конструктор & исключение и условие его выброса; второй тег~--- синоним первого \\
\texttt{@author}, \texttt{@version} & класс, интерфейс & автор и версия класса \\
\texttt{@since} версия & любой элемент & с какой версии элемент существует \\
\texttt{@see} ссылка & любой элемент & раздел «См. также» \\
\texttt{@deprecated} причина & любой элемент & элемент устарел; чем его заменить \\
\bottomrule
\end{tabular}
\end{table}

Тег \texttt{@param} пишется по одному на каждый параметр, в порядке сигнатуры; для обобщённых типов имя параметра типа берётся в угловые скобки. Тег \texttt{@return} у метода ровно один и для методов \texttt{void} не пишется. Теги \texttt{@throws} и \texttt{@exception}~--- одно и то же: исторически первым появился второй, затем его переименовали, а старое имя оставили для совместимости. У \texttt{@see} четыре формы: имя класса, имя класса с методом через решётку, внешняя ссылка и печатный источник текстом в кавычках; дописать пояснение в ту же строку нельзя. Тег \texttt{@deprecated} работает в паре с аннотацией \texttt{@Deprecated}: аннотация нужна компилятору, чтобы он выдал предупреждение, а тег~--- человеку, чтобы объяснить, чем заменить.

Строчных дескрипторов семь. Ссылки \texttt{\{@link Класс\#метод\}} и \texttt{\{@linkplain $\ldots$\}} отличаются только шрифтом~--- моноширинным и обычным; так же различаются \texttt{\{@code\}} и \texttt{\{@literal\}}, выводящие текст как есть. Тег \texttt{\{@value \#КОНСТАНТА\}} подставляет значение константы, \texttt{\{@inheritDoc\}} вставляет документацию переопределяемого метода, \texttt{\{@summary\}} задаёт границу краткого описания. Без \texttt{\{@code\}} строка \texttt{List<String>} превратится в \texttt{List}: браузер съест угловые скобки как неизвестный тег. А \texttt{\{@value\}} держит документацию в согласии с кодом: написав «максимум равен 500 рублей» текстом, вы обязаны помнить об этой фразе при каждой правке константы. Верх листинга~\ref{lst:10-2} показывает эти теги, низ~--- \texttt{\{@inheritDoc\}}.

\begin{listing}[H]
\begin{minted}{java}
/**
 * Возвращает список названий книг.
 *
 * <p>Тип результата - {@code List<String>}, никогда не
 * {@code null}. Условие: {@literal 0 < размер < 1000}.
 * Максимальный штраф: {@value FineCalculator#MAX_FINE} руб.
 */
List<String> titles();

// в классе-наследнике, у переопределённого метода:
/**
 * {@inheritDoc}
 *
 * <p>Реализация обращается к базе и не кэширует результат.
 */
@Override
public List<Student> findAll() {
    return studentRepository.findAll();
}
\end{minted}
\caption{Строчные дескрипторы и наследование документации}
\label{lst:10-2}
\end{listing}

Если Javadoc у переопределяющего метода отсутствует вообще, документация наследуется от родителя автоматически; \texttt{\{@inheritDoc\}} нужен, когда вы хотите взять описание родителя \emph{и добавить своё}.

\section{Генерация документации}

Утилита \texttt{javadoc} лежит в каталоге \texttt{bin} вашего JDK рядом с \texttt{java} и \texttt{javac}. Простейший вызов принимает один файл, реальный проект обрабатывают по пакетам (листинг~\ref{lst:10-3}); в PowerShell обратный слэш переноса не работает, поэтому команду там набирают одной строкой.

\begin{listing}[H]
\begin{minted}{bash}
javadoc -d docs -sourcepath src/main/java \
        -subpackages ru.fa.library \
        -encoding UTF-8 -charset UTF-8 -docencoding UTF-8 \
        -author -version -windowtitle "Онлайн-библиотека"
\end{minted}
\caption{Вызов утилиты javadoc из командной строки}
\label{lst:10-3}
\end{listing}

Опция \texttt{-d} задаёт каталог вывода, \texttt{-sourcepath}~--- где искать исходники, \texttt{-subpackages}~--- какой пакет обработать вместе с вложенными. Кодировок три: \texttt{-encoding} относится к исходным файлам, \texttt{-charset} и \texttt{-docencoding}~--- к сгенерированному HTML; без них русский текст в браузере превратится в мусор. Опция \texttt{-private} включает приватные члены, \texttt{-Xdoclint:none} отключает строгую проверку. Теги \texttt{@author} и \texttt{@version} по умолчанию игнорируются: студенты часто пишут \texttt{@author}, не видят его в HTML и решают, что тег не работает,~--- нужно просто добавить одноимённые флаги.

\begin{vrezka}
\textbf{Важно: doclint делит замечания на два класса.} Отсутствующий \texttt{@param} или \texttt{@return}~--- это \emph{предупреждение}: в консоли появится \texttt{warning: no @param for bookCount}, но документация сгенерируется, а код возврата останется нулевым. Битая ссылка в \texttt{\{@link\}} или \texttt{@see}, невалидный HTML и синтаксическая ошибка в комментарии~--- это \emph{ошибки}: утилита печатает \texttt{error: reference not found} и завершается с ненулевым кодом, HTML при этом не создаётся. Каталог вывода не очищается, там остаётся документация от прошлого удачного запуска, поэтому ориентируйтесь на консоль и код возврата, а не на наличие файлов.
\end{vrezka}

Руками команду \texttt{javadoc} в реальном проекте почти никто не вызывает~--- это делает система сборки. Разница как между стиркой в тазу и стиральной машиной: результат тот же, но во втором случае вы не стоите над ним. Настройка занимает один блок в файле \texttt{pom.xml} (листинг~\ref{lst:10-4}).

\begin{listing}[H]
\begin{minted}{xml}
<plugin>
    <groupId>org.apache.maven.plugins</groupId>
    <artifactId>maven-javadoc-plugin</artifactId>
    <version>3.11.2</version>
    <configuration>
        <encoding>UTF-8</encoding>
        <charset>UTF-8</charset>
        <docencoding>UTF-8</docencoding>
        <author>true</author>
        <version>true</version>
        <doclint>none</doclint>
    </configuration>
    <executions>
        <execution>
            <id>attach-javadocs</id>
            <goals><goal>jar</goal></goals>
        </execution>
    </executions>
</plugin>
\end{minted}
\caption{Подключение maven-javadoc-plugin}
\label{lst:10-4}
\end{listing}

Команда \texttt{mvn javadoc:javadoc} генерирует HTML, \texttt{mvn javadoc:jar} упаковывает документацию в архив, а \texttt{mvn package} собирает оба. Результат попадает в каталог \texttt{target/site/apidocs} или \texttt{target/reports/apidocs}; точный путь плагин печатает в консоли. Архив с суффиксом \texttt{-javadoc} нужен при публикации библиотеки: среда разработки скачивает его и показывает вашу документацию во всплывающих подсказках. Отдельный класс документируется своим комментарием, а весь пакет~--- файлом \texttt{package-info.java}, где нет ничего, кроме комментария и объявления пакета; его текст попадает на страницу \texttt{package-summary.html}.

\section{Зачем нужны автоматические тесты}

Пока проект маленький, вы тестируете его руками: запустили, нажали пару кнопок, посмотрели в консоль. Это \term{ручное тестирование}, и работает оно до того момента, когда экранов становится больше пяти: если полная проверка занимает 20~минут, а вы правите код 10~раз в день, на неё уйдёт больше трёх часов. Вы не станете этого делать~--- проверите только то, что правили, и сломается именно другое место. \term{Автоматизированное тестирование}~--- это программы, которые проверяют вашу программу: написали один раз, запускаете тысячу раз.

Подходов к проверке два: при подходе \term{чёрного ящика} (\eng{black box}) известен только внешний контракт~--- какие входы, какие ожидаемые выходы; при подходе \term{белого ящика} (\eng{white box}) известна внутренняя структура, и тесты пишутся так, чтобы пройти по всем ветвлениям и границам. Модульные тесты ближе ко второму. Сами тесты делятся на три уровня: \term{модульные} (\eng{unit}) проверяют один класс в изоляции, \term{интеграционные}~--- что несколько компонентов работают вместе, \term{сквозные} (\eng{end-to-end}) прогоняют полный сценарий через всё приложение. Соотношение принято изображать пирамидой (рис.~\ref{fig:10-1}).

\begin{figure}[htbp]\centering
\begin{tikzpicture}[font=\footnotesize]
\draw[thick] (0,0) -- (7,0) -- (3.5,4.2) -- cycle;
\draw (1.17,1.4) -- (5.83,1.4);
\draw (2.33,2.8) -- (4.67,2.8);
\node at (3.5,0.6) {Модульные (unit)};
\node at (3.5,1.95) {Интеграционные};
\node at (3.5,3.15) {\scriptsize E2E};
\node[anchor=west,align=left] at (7.7,0.6)
  {сотни тестов,\\миллисекунды на прогон};
\node[anchor=west,align=left] at (7.7,1.95)
  {десятки тестов,\\секунды на прогон};
\node[anchor=west,align=left] at (7.7,3.15)
  {единицы,\\минуты на прогон};
\end{tikzpicture}
\caption{Пирамида тестирования}\label{fig:10-1}
\end{figure}

Логика пирамиды проста: чем ниже уровень, тем дешевле тест и тем точнее он указывает на место поломки. Упавший модульный тест говорит: «метод \texttt{calculate} вернул 505 вместо 500»; упавший сквозной~--- «страница не открылась», и дальше ищите сами. Аналогия из автосервиса: каждую деталь проверяют на стенде, потом собранный узел и только затем выгоняют машину на тест-драйв. Перевёрнутую пирамиду называют антипаттерном \term{мороженое} (\eng{ice cream cone}): такая система тестируется часами и ломается постоянно.

Хороший модульный тест обязан быть \emph{быстрым} (тест, который лезет в реальную базу или в сеть, уже не модульный), \emph{независимым} (не зависит от порядка запуска~--- JUnit сознательно не гарантирует порядок методов) и \emph{повторяемым} (даёт один и тот же результат и на вашей машине, и на сервере сборки; тест, зависящий от текущей даты, сломается). К этим трём добавляют ещё два, получая аббревиатуру \eng{FIRST}: \eng{Fast, Independent, Repeatable, Self-validating} (тест сам решает, прошёл он или нет) и \eng{Timely} (пишется вместе с кодом). Любой тест состоит из трёх шагов~--- схема \term{AAA}: \eng{Arrange} (подготовка данных), \eng{Act} (вызов ровно одного тестируемого метода), \eng{Assert} (сверка результата с ожидаемым); в BDD ту же схему называют «дано~--- когда~--- тогда», а листинг~\ref{lst:10-5} показывает её в чистом виде. Разделяйте три части пустой строкой: тест, где подготовка перемешана с проверками, читать невозможно.

\begin{listing}[H]
\begin{minted}{java}
@Test
void calculate_returnsRateTimesDays() {
    FineCalculator calc = new FineCalculator();   // Arrange

    double fine = calc.calculate(10);             // Act

    assertEquals(50.0, fine, 0.001);              // Assert
}
\end{minted}
\caption{Тест, построенный по схеме AAA}
\label{lst:10-5}
\end{listing}

\section{Фреймворк JUnit 5}

\term{JUnit}~--- самый распространённый фреймворк модульного тестирования для Java. Пятая версия переписана заново и разделена на три подпроекта. \term{JUnit Platform}~--- фундамент: запускает тесты и связывает их со средой разработки, Maven и командной строкой; сама тесты не выполняет, а определяет интерфейс \texttt{TestEngine}, через который к ней подключаются «движки». \term{JUnit Jupiter}~--- модель программирования: аннотации \texttt{@Test}, \texttt{@BeforeEach}, \texttt{@DisplayName}, класс \texttt{Assertions} и механизм расширений; собран из артефактов \texttt{junit-jupiter-api}, \texttt{junit-jupiter-engine} и \texttt{junit-jupiter-params}. \term{JUnit Vintage}~--- движок обратной совместимости для тестов JUnit~3 и~4. Аналогия: платформа~--- розетка в стене, Jupiter и Vintage~--- вилки разных стандартов, ваши тесты~--- приборы.

В обычном Maven-проекте достаточно одной зависимости-агрегата и плагина Surefire (листинг~\ref{lst:10-6}). Обратите внимание на область видимости \texttt{test}: библиотека нужна только при компиляции и запуске тестов, в итоговый архив приложения она не попадёт. В Spring Boot ещё проще: стартер \texttt{spring-boot-starter-test} приносит JUnit~5, Mockito, AssertJ и модуль Spring Test одной строкой.

\begin{listing}[H]
\begin{minted}{xml}
<dependencies>
    <dependency>
        <groupId>org.junit.jupiter</groupId>
        <!-- зонтичный артефакт: api + engine + params -->
        <artifactId>junit-jupiter</artifactId>
        <version>5.11.4</version>
        <scope>test</scope>
    </dependency>
</dependencies>

<build><plugins>
    <plugin>
        <groupId>org.apache.maven.plugins</groupId>
        <artifactId>maven-surefire-plugin</artifactId>
        <version>3.5.2</version>
    </plugin>
</plugins></build>
\end{minted}
\caption{Подключение JUnit 5 к Maven-проекту}
\label{lst:10-6}
\end{listing}

Где живут тесты, показано на рис.~\ref{fig:10-2}: в каталоге \texttt{src/test/java}, в том же пакете, что и тестируемый класс,~--- тогда тест видит его члены с пакетной видимостью. Имя тестового класса заканчивается на \texttt{Test} или \texttt{Tests}: Surefire ищет именно такие, а классы с суффиксом \texttt{IT} запускает плагин Failsafe.

\begin{figure}[htbp]\centering
\begin{forest} hierarchy
[project/
  [pom.xml]
  [src/
    [main/java/ru/fa/library/service/FineCalculator.java]
    [test/java/ru/fa/library/service/FineCalculatorTest.java]
  ]
  [target/]
]
\end{forest}
\caption{Размещение кода и тестов в Maven-проекте}\label{fig:10-2}
\end{figure}

Команда \texttt{mvn test} запускает все тесты; ключ \texttt{-Dtest} с именем класса ограничивает прогон одним классом, а имя класса с решёткой и именем метода~--- одним методом.

\subsection{Тестовые классы, методы и жизненный цикл}

\term{Тестовый класс}~--- обычный класс с хотя бы одним тестовым методом; он не должен быть абстрактным и должен иметь единственный конструктор. \term{Тестовый метод} помечен \texttt{@Test} или родственной аннотацией: он не может быть статическим или приватным, возвращает \texttt{void} и не принимает параметров, кроме тех, что внедряет сам JUnit. Важное отличие от JUnit~4: тестовые классы и методы \emph{не обязаны} быть публичными. Имя теста должно читаться как утверждение о поведении (схемы \texttt{метод\_ожидаемоеПоведение} или \texttt{should\_поведение\_when\_условие}), а поскольку оно ограничено синтаксисом Java, существует \texttt{@DisplayName}: она задаёт имя для человека, и отчёт становится живым описанием поведения системы.

Метод с аннотацией \texttt{@BeforeEach} выполняется перед каждым тестовым методом класса, метод с \texttt{@AfterEach}~--- после каждого, даже если тест упал. Бытовая аналогия: разложить доску, нож и продукты перед каждым блюдом и вымыть их после~--- не «один раз в начале смены», чтобы вкус предыдущего блюда не попал в следующее. Разберём это на классе с изменяемым состоянием. Пусть \texttt{BookShelf}~--- книжная полка: внутри список \texttt{ArrayList} с названиями, метод \texttt{add} отвергает пустое название исключением \texttt{IllegalArgumentException}, \texttt{remove} убирает книгу, \texttt{size} даёт количество, \texttt{titles}~--- неизменяемую копию списка. Тест к нему приведён в листинге~\ref{lst:10-7}. Ключевой момент: JUnit создаёт \emph{новый экземпляр} тестового класса для каждого тестового метода, поэтому поле \texttt{shelf} в первом тесте и во втором~--- два разных объекта. Именно так обеспечивается независимость, и второй тест это проверяет.

\begin{listing}[H]
\begin{minted}{java}
@DisplayName("Книжная полка")
class BookShelfTest {

    private BookShelf shelf;

    @BeforeEach
    void setUp() {          // перед КАЖДЫМ тестом
        shelf = new BookShelf();
        shelf.add("Война и мир");
    }

    @AfterEach              // после каждого, даже если упал
    void tearDown() {
        System.out.println("книг на полке: " + shelf.size());
    }

    @Test
    @DisplayName("После добавления книги размер увеличивается")
    void add_increasesSize() {
        shelf.add("Мастер и Маргарита");

        assertEquals(2, shelf.size(), "Должно быть две книги");
    }

    @Test
    @DisplayName("Каждый тест получает свежую полку")
    void eachTestGetsFreshShelf() {
        assertEquals(1, shelf.size());
        assertFalse(shelf.titles().contains("Мастер и Маргарита"));
    }
}
\end{minted}
\caption{Тест с аннотациями жизненного цикла}
\label{lst:10-7}
\end{listing}

Иногда подготовка слишком дорога, чтобы повторять её перед каждым тестом: поднять базу, установить соединение. Метод с аннотацией \texttt{@BeforeAll} выполняется один раз перед всеми тестами класса, с \texttt{@AfterAll}~--- один раз после всех; экземпляра класса на этот момент ещё не существует, поэтому по умолчанию такие методы обязаны быть статическими. Правило: в \texttt{@BeforeAll}~--- то, что не меняется от теста к тесту, всё изменяемое состояние~--- в \texttt{@BeforeEach}. Полный порядок выполнения показан на рис.~\ref{fig:10-3}; обратите внимание на середину схемы: конструктор вызывается дважды.

\begin{figure}[htbp]\centering
\begin{tikzpicture}[font=\footnotesize,
  bx/.style={draw,rounded corners=2pt,align=center,text width=25mm,
             minimum height=9mm,inner sep=2pt}]
\node[bx,text width=42mm](ba)at(6.5,5.1){\texttt{@BeforeAll}~--- один раз};
\node[bx](c1)at(1.6,3.4){конструктор\\\scriptsize экземпляр 1};
\node[bx](b1)at(4.9,3.4){\texttt{@BeforeEach}};
\node[bx](t1)at(8.2,3.4){\texttt{@Test} test1};
\node[bx](a1)at(11.5,3.4){\texttt{@AfterEach}};
\node[bx](c2)at(1.6,1.7){конструктор\\\scriptsize экземпляр 2};
\node[bx](b2)at(4.9,1.7){\texttt{@BeforeEach}};
\node[bx](t2)at(8.2,1.7){\texttt{@Test} test2};
\node[bx](a2)at(11.5,1.7){\texttt{@AfterEach}};
\node[bx,text width=42mm](aa)at(6.5,0){\texttt{@AfterAll}~--- один раз};
\draw[flowarrow](ba.south)--++(0,-0.3)-|(c1.north);
\draw[flowarrow](c1)--(b1);\draw[flowarrow](b1)--(t1);
\draw[flowarrow](t1)--(a1);
\draw[flowarrow](a1.south)--++(0,-0.3)-|(c2.north);
\draw[flowarrow](c2)--(b2);\draw[flowarrow](b2)--(t2);
\draw[flowarrow](t2)--(a2);
\draw[flowarrow](a2.south)--++(0,-0.3)-|(aa.north);
\end{tikzpicture}
\caption{Жизненный цикл тестового класса с двумя тестами}\label{fig:10-3}
\end{figure}

Такое поведение включено по умолчанию~--- это режим \texttt{PER\_METHOD} аннотации \texttt{@TestInstance}; режим \texttt{PER\_CLASS} создаёт один экземпляр на все тесты и разрешает нестатический \texttt{@BeforeAll}, но платой становится общее состояние. Остальные аннотации сведены в табл.~\ref{tab:10-2}. Три стоит пояснить: \texttt{@Disabled} временно отключает тест~--- обязательно указывайте причину, иначе через месяц никто не вспомнит, почему он выключен; \texttt{@Tag} помечает тесты категориями, и тогда \texttt{mvn test -Dgroups=fast} запустит только помеченные; \texttt{@Nested} группирует тесты во вложенных классах, причём вложенный класс обязан быть нестатическим внутренним.

\begin{table}[htbp]
\caption{Основные аннотации JUnit 5}
\label{tab:10-2}
\centering
\begin{tabular}{L{4.6cm}L{10.4cm}}
\toprule
Аннотация & Что делает \\
\midrule
\texttt{@Test}, \texttt{@DisplayName} & тестовый метод; его имя для человека в отчёте \\
\texttt{@BeforeEach}, \texttt{@AfterEach} & выполнить перед каждым тестом и после каждого \\
\texttt{@BeforeAll}, \texttt{@AfterAll} & выполнить один раз до и после всех тестов класса \\
\texttt{@Disabled}, \texttt{@Tag} & отключить тест; пометить категорией \\
\texttt{@Nested}, \texttt{@ParameterizedTest} & вложенная группа; тест на наборах данных \\
\texttt{@RepeatedTest}, \texttt{@TestFactory} & повторить тест N раз; фабрика динамических тестов \\
\texttt{@Timeout}, \texttt{@ExtendWith} & ограничить время; подключить расширение \\
\bottomrule
\end{tabular}
\end{table}

В старом коде встретятся аннотации JUnit~4: \texttt{@Before} и \texttt{@After} соответствуют \texttt{@BeforeEach} и \texttt{@AfterEach}, \texttt{@Ignore}~--- \texttt{@Disabled}, \texttt{@RunWith}~--- \texttt{@ExtendWith}. Смешивать две версии в одном классе нельзя.

\section{Утверждения и проверка исключений}

\term{Утверждение} (\eng{assertion})~--- проверка условия внутри теста: условие выполнено~--- тест продолжается, не выполнено~--- метод бросает ошибку и тест помечается как провалившийся. Аналогия: утверждение~--- калибр у контролёра; деталь либо проходит в шаблон, либо нет, оценки «почти подошла» не бывает.

Здесь нужно расставить точки над «i», потому что этот вопрос путает всех. В JUnit~4 класс утверждений назывался \texttt{Assert} и лежал в пакете \texttt{org.junit}; в JUnit~5 он называется \texttt{Assertions} и лежит в \texttt{org.junit.jupiter.api}. Изменилось и содержание: появились \texttt{assertThrows} и \texttt{assertAll}, а сообщение об ошибке переехало из первого аргумента в последний. Если в коде импортируется \texttt{org.junit.Assert}~--- вы пишете тест на JUnit~4, даже если думаете иначе.

Основные методы собраны в листинге~\ref{lst:10-8}. Держите в голове два правила. Первое: порядок аргументов критичен~--- сначала \emph{ожидаемое} значение, потом \emph{фактическое}; перепутаете~--- тест будет работать, но сообщение об ошибке введёт вас в заблуждение. Второе: \texttt{assertEquals} вызывает метод \texttt{equals}, а \texttt{assertSame} сравнивает ссылки оператором \texttt{==}. У каждого метода есть перегрузка с сообщением. Отдельная тема, на которой спотыкаются все,~--- вещественные числа: \texttt{double} и \texttt{float} хранятся приближённо (см. главу~\ref{ch:01}), поэтому третьим аргументом \texttt{assertEquals} принимает \term{дельту}~--- максимально допустимое расхождение. Для денег её берут порядка копейки, для физических расчётов~--- по требуемой точности.

\begin{listing}[H]
\begin{minted}{java}
assertEquals(4, 2 + 2);      // ожидаемое, потом фактическое
assertTrue(4 > 3);
assertFalse("текст".isBlank());
assertNull(nothing);
assertNotNull("текст");
assertSame(first, second);   // идентичность ссылок (==)
assertArrayEquals(new int[]{1, 2}, new int[]{1, 2});
fail("сюда исполнение попадать не должно");

@Test
void floatingPointIsTricky() {
    double sum = 0.1 + 0.2;

    // assertEquals(0.3, sum);   упал бы: 0.30000000000000004
    assertEquals(0.3, sum, 0.000001);   // с допуском (дельтой)
}
\end{minted}
\caption{Методы Assertions, дельта и групповая проверка}
\label{lst:10-8}
\end{listing}

Отдельного упоминания заслуживает \texttt{assertAll}. Обычное утверждение прерывает тест сразу, поэтому вы узнаёте только о первой сломавшейся проверке; \texttt{assertAll} выполняет все переданные ему лямбды и показывает все ошибки сразу~--- его применяют для независимых проверок одного объекта, и в действии он виден в листинге~\ref{lst:10-9}.

\subsection{Тестирование исключений}

Метод обязан не только возвращать правильный результат, но и правильно \emph{отказываться} работать. Аналогия: предохранитель проверяют не тем, что он пропускает ток, а тем, что при перегрузке он сгорает; не сгорел~--- значит бракованный, даже если лампочка горит.

В JUnit~4 для этого был атрибут \texttt{@Test(expected = $\ldots$)} с двумя изъянами: непонятно, какая именно строка бросила исключение, и нельзя проверить текст сообщения. JUnit~5 решает задачу методом \texttt{assertThrows}: он принимает класс ожидаемого исключения и лямбду с кодом, который должен упасть,~--- тест проходит, только если брошено исключение нужного типа. Главное преимущество в том, что \texttt{assertThrows} \emph{возвращает пойманное исключение}, и его можно исследовать дальше: сообщение, код ошибки, вложенную причину (листинг~\ref{lst:10-9}). Граничные значения~--- ноль, минус единица, максимум плюс-минус единица~--- проверяют первыми: именно там живёт большинство ошибок.

\begin{listing}[H]
\begin{minted}{java}
@DisplayName("Калькулятор штрафов: границы и отказы")
class FineCalculatorExceptionTest {

    private FineCalculator calc;

    @BeforeEach
    void setUp() { calc = new FineCalculator(); }

    @Test
    @DisplayName("Сообщение об ошибке содержит само значение")
    void negativeDays_throwsWithMessage() {
        IllegalArgumentException e = assertThrows(
                IllegalArgumentException.class,
                () -> calc.calculate(-5));

        assertAll(
            () -> assertTrue(e.getMessage().contains("отрицат")),
            () -> assertTrue(e.getMessage().contains("-5")));
    }
}
\end{minted}
\caption{Проверка отказов и граничных значений}
\label{lst:10-9}
\end{listing}

Два уточнения. Метод \texttt{assertThrows} принимает и наследников указанного класса: если код бросит \texttt{NumberFormatException} (наследник \texttt{IllegalArgumentException}), тест с ожиданием родителя пройдёт; когда нужен ровно этот класс, есть \texttt{assertThrowsExactly}. И строго говоря, любое непойманное исключение и так провалит тест~--- метод \texttt{assertDoesNotThrow} нужен ради читаемости и заодно возвращает результат вызова.

\section{Наборы данных: параметризованные и динамические тесты}

Посмотрите на четыре тестовых метода, различающиеся только числом дней и ожидаемой суммой, и почувствуйте раздражение: логика одна, данные разные. Это как отдельный рецепт для каждой партии котлет~--- разумнее написать один рецепт и подавать в него разные наборы продуктов.

\term{Параметризованный тест} выполняется столько раз, сколько наборов данных подал источник; нужен артефакт \texttt{junit-jupiter-params} из зонтичного агрегата. Шаблон имени прогона задаётся атрибутом \texttt{name} самой аннотации \texttt{@ParameterizedTest}~--- ни у одного источника такого атрибута нет; в шаблоне доступны плейсхолдеры \texttt{\{0\}} и \texttt{\{1\}} для аргументов и \texttt{\{index\}} для номера прогона. Источников шесть: \texttt{@ValueSource}~--- массив литералов одного типа в единственный параметр; \texttt{@NullSource}, \texttt{@EmptySource} и объединяющая их \texttt{@NullAndEmptySource}~--- \texttt{null} и пустое значение; \texttt{@CsvSource}~--- строки CSV с несколькими параметрами на прогон, а \texttt{@CsvFileSource} читает их из файла ресурсов; \texttt{@EnumSource}~--- значения перечисления; \texttt{@MethodSource}~--- результат статического метода. Первые три показаны в листинге~\ref{lst:10-10}, и второй тест демонстрирует важную деталь: \texttt{@ValueSource} не умеет передавать \texttt{null}, потому что аннотации Java не принимают \texttt{null} в массивах значений.

\begin{listing}[H]
\begin{minted}{java}
@ParameterizedTest(name = "просрочка {0} дней - ошибки нет")
@ValueSource(ints = {0, 1, 5, 100, 1000})
void calculate_acceptsNonNegativeDays(int days) {
    assertDoesNotThrow(() -> calc.calculate(days));
}

@ParameterizedTest(name = "название «{0}» отвергается")
@NullAndEmptySource                  // прогон с null и с ""
@ValueSource(strings = {"   ", "\t"})
void add_rejectsBlankTitles(String title) {
    assertThrows(IllegalArgumentException.class,
                 () -> new BookShelf().add(title));
}

@ParameterizedTest(name = "{0} дней -> {1} руб.")
@CsvSource({"0, 0.0", "1, 5.0", "100, 500.0", "999, 500.0"})
void calculate_matchesTable(int days, double expected) {
    assertEquals(expected, calc.calculate(days), 0.001);
}
\end{minted}
\caption{Источники @ValueSource, @NullAndEmptySource и @CsvSource}
\label{lst:10-10}
\end{listing}

Когда данные~--- не литералы, а объекты, или когда к каждому набору хочется приложить описание, берут \texttt{@MethodSource}: статический метод в том же классе возвращает поток объектов \texttt{Arguments} (листинг~\ref{lst:10-11}). Статическим он обязан быть потому, что источник данных вычисляется до создания экземпляра тестового класса.

\begin{listing}[H]
\begin{minted}{java}
@ParameterizedTest(name = "{2}")
@MethodSource("fineCases")
void calculate_worksForAllCases(int days, double expected,
                                String description) {
    assertEquals(expected, calc.calculate(days), 0.001);
}

/** Источник данных: дни, ожидаемый штраф, пояснение. */
static Stream<Arguments> fineCases() {
    return Stream.of(
        Arguments.of(0,   0.0,   "без просрочки штрафа нет"),
        Arguments.of(99,  495.0, "до максимума считаем линейно"),
        Arguments.of(100, 500.0, "ровно максимум"),
        Arguments.of(101, 500.0, "выше максимума не растёт"));
}
\end{minted}
\caption{Источник @MethodSource с описанием каждого случая}
\label{lst:10-11}
\end{listing}

Ещё две аннотации решают смежные задачи. \texttt{@RepeatedTest} выполняет один и тот же тест заданное число раз~--- это нужно коду со случайностью, конкурентностью или кэшированием, где ошибка проявляется не всегда; в шаблоне имени доступны плейсхолдеры \texttt{\{currentRepetition\}} и \texttt{\{totalRepetitions\}}. Помните ограничение: если ошибка случается в одном прогоне из тысячи, десять повторов её не поймают~--- это инструмент диагностики, а не гарантия. \texttt{@TestFactory} строит \term{динамические тесты}: метод возвращает не \texttt{void}, а поток объектов \texttt{DynamicTest} со своими именами-строками, поэтому набор проверок может зависеть от данных, прочитанных во время выполнения; если же данные известны при компиляции, параметризованный тест проще.

\section{Mockito: тестирование в изоляции}

Модульный тест обязан проверять один класс. Но реальный сервис почти всегда зависит от других объектов: репозитория, отправителя писем, HTTP-клиента. Возьмём сервис выдачи книг из листинга~\ref{lst:10-12}: он получает книгу из хранилища, проверяет доступность, сохраняет изменение и уведомляет читателя. Чтобы протестировать метод \texttt{lend} «по-настоящему», пришлось бы поднять базу и почтовый сервер: тест стал бы медленным и~--- самое неприятное~--- реально слал бы письма.

\begin{listing}[H]
\begin{minted}{java}
package ru.fa.library.service;

public interface BookRepository {
    Optional<Book> findById(Long id);
    Book save(Book book);
}

public interface NotificationSender {
    void send(String email, String text);
}

/** Сервис выдачи книг читателям. */
public class LoanService {

    private final BookRepository repository;
    private final NotificationSender sender;

    // конструктор принимает обе зависимости и сохраняет их

    /**
     * Выдаёт книгу читателю и отправляет уведомление.
     *
     * @throws IllegalArgumentException если книги нет
     * @throws IllegalStateException    если она уже выдана
     */
    public void lend(Long id, String email) {
        Book book = repository.findById(id).orElseThrow(
            () -> new IllegalArgumentException("Нет: " + id));
        if (!book.isAvailable()) {
            throw new IllegalStateException(
                    "Книга уже выдана: " + book.getTitle());
        }
        book.setAvailable(false);
        repository.save(book);
        sender.send(email, "Вы получили «" + book.getTitle() + "»");
    }
}
\end{minted}
\caption{Сервис с двумя зависимостями}
\label{lst:10-12}
\end{listing}

Класс \texttt{Book} прост: конструктор принимает идентификатор, название и признак доступности, а наружу отдаёт геттеры и сеттер признака. Аналогия к решению: на тренинге кассиров банкомат заменяют картонным макетом~--- он выдаёт «купюры» по команде, но денег не тратит; проверяют при этом не банкомат, а кассира. \term{Mockito}~--- библиотека, создающая такие макеты для Java-объектов.

Терминология вокруг тестовых дублёров сложилась богатая. \eng{Dummy}~--- пустышка, нужная лишь чтобы заполнить параметр. \eng{Stub} (заглушка) возвращает заранее заданные ответы, и мы проверяем состояние результата. \eng{Mock} (имитация) добавляет ожидания по вызовам: проверяется поведение~--- что и сколько раз вызвали. \eng{Spy} (шпион)~--- обёртка над настоящим объектом. \eng{Fake}~--- упрощённая рабочая реализация, например репозиторий на \texttt{HashMap}. Практическое различие заглушки и мока укладывается в одну фразу: \emph{заглушка отвечает на вопросы, мок проверяет вопросы},~--- а в Mockito это один объект, роль определяется тем, как вы его используете.

Подключается библиотека зависимостью \texttt{mockito-junit-jupiter} с областью видимости \texttt{test}: она тянет за собой ядро и добавляет расширение для JUnit~5, а в Spring Boot входит в тестовый стартер. Создать макет можно программно вызовом \texttt{Mockito.mock}, но предпочтительнее аннотации~--- каркас теста и удачный сценарий приведены в листинге~\ref{lst:10-13}.

\begin{listing}[H]
\begin{minted}{java}
import static org.mockito.ArgumentMatchers.*;  // any, eq
import static org.mockito.Mockito.*;           // when, verify

@ExtendWith(MockitoExtension.class)
@DisplayName("Сервис выдачи книг")
class LoanServiceTest {

    @Mock  private BookRepository repository;
    @Mock  private NotificationSender sender;

    @InjectMocks
    private LoanService service;   // моки уйдут в конструктор

    @Test
    @DisplayName("Книга выдаётся, сохраняется, уходит письмо")
    void lend_marksBookUnavailableAndNotifies() {
        Book book = new Book(1L, "Война и мир", true);
        when(repository.findById(1L))
                .thenReturn(Optional.of(book));

        service.lend(1L, "ivan@example.com");

        assertFalse(book.isAvailable());
        verify(repository).save(book);
        verify(sender).send(eq("ivan@example.com"), anyString());
    }
}
\end{minted}
\caption{Настройка макетов и проверка вызовов}
\label{lst:10-13}
\end{listing}

Расширение \texttt{MockitoExtension} перед каждым тестом создаёт свежие макеты и внедряет их в поле, помеченное \texttt{@InjectMocks}: настройки и записанные вызовы не должны переноситься из теста в тест. Новосозданный макет возвращает «пустое» значение из любого метода~--- \texttt{null} для объектов, ноль для чисел, \texttt{false} для логического типа, пустой \texttt{Optional} и пустую коллекцию. Конструкция \texttt{when($\ldots$).thenReturn($\ldots$)} читается почти как русский текст: «когда у репозитория спросят книгу номер один, тогда вернуть вот эту». Метод \texttt{thenThrow} заставляет макет бросить исключение (для методов \texttt{void} порядок обратный: \texttt{doThrow($\ldots$).when($\ldots$)}). Число вызовов задаётся вторым аргументом \texttt{verify}~--- \texttt{times(1)} по умолчанию, а также \texttt{never()} и \texttt{atLeastOnce()}; проверка «больше ничего не было»~--- \texttt{verifyNoInteractions}.

\begin{oshibka}
\textbf{Частая ошибка: смешивание матчеров и обычных значений.} \term{Матчеры аргументов}~--- \texttt{any()}, \texttt{anyString()}, \texttt{eq(значение)}, \texttt{argThat(предикат)}~--- позволяют не указывать точные значения. Железное правило: если хотя бы один аргумент задан матчером, остальные тоже обязаны быть матчерами. Вызов \texttt{verify(sender).send(адрес, anyString())} с «сырым» адресом приведёт к исключению \texttt{InvalidUseOfMatchersException}; правильно обернуть точное значение в \texttt{eq($\ldots$)}.
\end{oshibka}

Метод \texttt{verify} отвечает на вопрос «вызвали ли», а на вопрос «с какими именно аргументами» отвечает \texttt{ArgumentCaptor}~--- он незаменим, когда аргумент собран внутри тестируемого метода. Захватчик объявляют как \texttt{ArgumentCaptor<String>}, создают вызовом \texttt{ArgumentCaptor.forClass}, подставляют \texttt{captor.capture()} на место аргумента внутри \texttt{verify} и затем читают \texttt{captor.getValue()}. Наконец, \term{шпион} оборачивает \emph{реальный} объект: методы работают как обычно, но вызовы записываются. Здесь есть ловушка: \texttt{when(spyList.get(5))} реально вызовет \texttt{get(5)} и бросит исключение, поэтому для шпиона используют форму \texttt{doReturn($\ldots$).when(spyList).get(5)}.

\section{Тестирование Spring-приложений и покрытие кода}

Вернёмся к приложению из главы~\ref{ch:07}: контроллер, сервис, репозиторий. Проверять его можно на трёх уровнях, и каждый стоит разного времени (табл.~\ref{tab:10-3}). Аналогия~--- приёмка квартиры: розетку проверяют тестером прямо в стене, потом весь электрощиток и только затем заселяются и включают всю технику разом.

\begin{table}[htbp]
\caption{Три уровня тестов Spring-приложения}
\label{tab:10-3}
\centering
\begin{tabular}{L{3.9cm}L{3.9cm}L{3.9cm}L{2.2cm}}
\toprule
Уровень & Аннотация & Что поднимается & Скорость \\
\midrule
Модульный & нет, только Mockito & ничего, обычные объекты & мс \\
Срез (\eng{slice}) & \texttt{@WebMvcTest}, \texttt{@DataJpaTest} & часть контекста Spring & сотни мс \\
Интеграционный & \texttt{@SpringBootTest} & весь контекст, при желании сервер & секунды \\
\bottomrule
\end{tabular}
\end{table}

Самый быстрый и самый нужный уровень~--- модульный тест сервиса, где Spring не участвует вообще: мы создаём объект и подсовываем ему макет репозитория, как в предыдущем разделе. Именно такими тестами покрывают основную массу бизнес-логики.

Следующий уровень~--- срез веб-слоя. Аннотация \texttt{@WebMvcTest} поднимает \emph{только} контроллеры, преобразователи JSON и обработчики ошибок; сервисы и репозитории в контекст не попадают, и их подменяют макетами. Инструмент \texttt{MockMvc} выполняет HTTP-запросы без сетевого соединения и без запуска сервера, вызывая диспетчер напрямую (листинг~\ref{lst:10-14}). Такой тест проверяет маршрутизацию адресов, коды ответа, формат JSON и разбор тела запроса, но принципиально не проверяет работу сервиса~--- тот заменён макетом.

\begin{listing}[H]
\begin{minted}{java}
@WebMvcTest(StudentRestController.class)
@DisplayName("REST-контроллер студентов")
class StudentRestControllerTest {

    @Autowired
    private MockMvc mockMvc;

    @MockitoBean          // подменяет бин сервиса макетом
    private StudentService studentService;

    @Test
    @DisplayName("GET /api/students отдаёт 200 и список")
    void getAll_returnsList() throws Exception {
        when(studentService.findAll()).thenReturn(
                List.of(new Student(1L, "Иван", "Петров")));

        mockMvc.perform(get("/api/students"))
                .andExpect(status().isOk())
                .andExpect(jsonPath("$[0].name").value("Иван"));
    }
}
\end{minted}
\caption{Срез веб-слоя: @WebMvcTest и MockMvc}
\label{lst:10-14}
\end{listing}

Аннотация \texttt{@MockitoBean} появилась в Spring Framework~6.2 (Spring Boot~3.4) на замену устаревшей \texttt{@MockBean}; в книгах более ранних лет встретится старое имя, смысл тот же~--- подменить бин в контексте макетом. Рядом живёт \texttt{@MockitoSpyBean}~--- обёртка-шпион вокруг настоящего бина.

Симметричный срез для слоя данных~--- \texttt{@DataJpaTest}: он поднимает только JPA-часть, подставляет встроенную базу в памяти и даёт помощник \texttt{TestEntityManager}, а каждый тест идёт в транзакции, которая после завершения откатывается. Так проверяют производные методы репозиториев вроде \texttt{findBySurname}: нигде, кроме теста, вы не убедитесь, что имя метода превратилось в правильный SQL. Стандартные \texttt{save} и \texttt{findById} тестировать не надо~--- вы проверяли бы чужую библиотеку.

Верхний уровень~--- \texttt{@SpringBootTest}: поднимается весь контекст приложения. Макетов здесь нет~--- работают настоящий сервис и настоящая сериализация JSON, а при значении \texttt{webEnvironment = RANDOM\_PORT} запрос идёт через сетевой стек к серверу на случайном свободном порту. Клиент \texttt{TestRestTemplate} уже знает адрес и порт, поэтому в запросах указывается только путь (листинг~\ref{lst:10-15}).

\begin{listing}[H]
\begin{minted}{java}
@SpringBootTest(
        webEnvironment = SpringBootTest.WebEnvironment.RANDOM_PORT)
@DisplayName("Сквозной сценарий: создание и чтение списка")
class StudentIntegrationTest {

    @Autowired
    private TestRestTemplate restTemplate;

    @Test
    void createdStudent_appearsInList() {
        Student request = new Student(null, "Анна", "Смирнова");

        ResponseEntity<Student> created = restTemplate
                .postForEntity("/api/students", request,
                               Student.class);

        assertEquals(HttpStatus.CREATED, created.getStatusCode());
        assertNotNull(created.getBody().id());

        ResponseEntity<Student[]> all = restTemplate
                .getForEntity("/api/students", Student[].class);

        assertTrue(Arrays.stream(all.getBody()).anyMatch(
                s -> "Анна".equals(s.name())));
    }
}
\end{minted}
\caption{Интеграционный тест всего приложения}
\label{lst:10-15}
\end{listing}

Атрибут \texttt{webEnvironment} принимает четыре значения: \texttt{MOCK} (по умолчанию~--- веб-окружение имитируется, реальный порт не занимается), \texttt{RANDOM\_PORT}, \texttt{DEFINED\_PORT} и \texttt{NONE}. Настройки для тестов держат в отдельном файле профиля и включают аннотацией \texttt{@ActiveProfiles("test")}. Правило выбора уровня таково: бизнес-логику покрывайте модульными тестами с макетами (их должно быть больше всего), контроллеры~--- срезом веб-слоя, нестандартные методы репозиториев~--- срезом слоя данных, а один-два ключевых сценария целиком~--- интеграционным тестом. Не наоборот: приложение, покрытое исключительно \texttt{@SpringBootTest}, собирается минутами и проваливается целиком при любой мелочи.

\subsection{Покрытие кода}

\term{Покрытие кода} (\eng{code coverage})~--- доля строк, ветвлений или методов, выполненных во время прогона тестов; метрика показывает, какие участки кода тесты ни разу не задели. Стандартный инструмент для Java~--- \term{JaCoCo}: он подключается агентом к виртуальной машине и строит отчёт в каталоге \texttt{target/site/jacoco}, где каждая строка подсвечена: зелёная выполнялась, красная ни разу, жёлтая~--- ветвление проверено частично. Основных метрик две: покрытие строк и покрытие ветвлений; вторая строже и полезнее.

\begin{oshibka}
\textbf{Частая ошибка: покрытие принимают за качество.} Покрытие говорит, что строка \emph{выполнилась}. Оно не говорит, что результат \emph{проверили}. Тест из одной строки \texttt{new FineCalculator().calculate(10);} без единого утверждения даёт стопроцентное покрытие метода и не проверяет ровным счётом ничего: строка выполнена, инструмент доволен, ошибка в формуле проедет мимо. Бытовая аналогия: покрытие~--- отчёт «прошёлся с пылесосом по всем комнатам». Прошёлся~--- да. Убрал ли~--- неизвестно: пылесос мог быть выключен.
\end{oshibka}

Отсюда практические выводы. Низкое покрытие~--- надёжный сигнал беды: если критичный сервис покрыт на 10~\%, там точно есть непроверенные сценарии. Высокое покрытие~--- не гарантия качества, а лишь необходимое условие, и погоня за 100~\% вредна: последние проценты приходятся на геттеры, недостижимые ветки и обработчики «этого не может быть», тесты на них дороги и мешают рефакторингу. Разумный ориентир~--- 70--80~\% по ветвлениям для бизнес-логики. Ориентируйтесь не на процент, а на вопрос: «Если я сломаю эту строчку, упадёт ли хоть один тест?» Если нет~--- покрытие фиктивное.

\praktikum{}

Понадобятся JDK~21 или новее, Maven~3.9 или новее и текстовый редактор. Задания выполняются по порядку: каждое опирается на код предыдущего. Задания~1--4 живут в проекте \texttt{library-tests}, задание~5 требует отдельного проекта на Spring Boot.

\subsection*{Задание 1. Проект и документация Javadoc}

Создайте вручную каталог \texttt{library-tests} с файлом \texttt{pom.xml} и двумя ветками исходников~--- \texttt{src/main/java} и \texttt{src/test/java},~--- заведя в каждой пакет \texttt{ru.fa.library.service}. В \texttt{pom.xml} объявите координаты \texttt{ru.fa}, \texttt{library-tests}, версию \texttt{1.0-SNAPSHOT}, свойства \texttt{maven.compiler.release} (значение~21) и \texttt{project.build.sourceEncoding} (значение \texttt{UTF-8}), подключите плагин \texttt{maven-compiler-plugin} и соберите проект командой \texttt{mvn clean compile}.

Положите в этот пакет класс \texttt{FineCalculator} из листинга~\ref{lst:10-1}, убрав все комментарии, и добавьте метод \texttt{calculateForBooks(int daysOverdue, int bookCount)}: он умножает результат \texttt{calculate} на количество книг и бросает \texttt{IllegalArgumentException}, если книг не больше нуля. Посмотрите на класс глазами человека, который видит его впервые: по сигнатуре не понять ни единиц результата, ни поведения при нуле. Напишите инструкцию: задокументируйте класс (описание существительным, ссылки \texttt{\{@link\}} на обе константы, теги \texttt{@author}, \texttt{@version}, \texttt{@since}), обе константы и оба метода (глагол в третьем лице, поведение на границах, \texttt{@param} на каждый параметр, \texttt{@return}, \texttt{@throws}, хотя бы один \texttt{\{@code\}} и один \texttt{\{@value\}}). Добавьте файл \texttt{package-info.java} с описанием пакета.

Сгенерируйте документацию командой из листинга~\ref{lst:10-3} и найдите в \texttt{docs/index.html} страницу класса, строку с автором и версией, подставленное значение вместо \texttt{\{@value\}} и страницу пакета. Затем два эксперимента, после каждого возвращая код: удалите тег \texttt{@param bookCount}~--- появится предупреждение, но документация обновится; замените \texttt{\{@link \#MAX\_FINE\}} на несуществующее имя~--- будет ошибка и ненулевой код возврата, а с флагом \texttt{-Xdoclint:none} документация соберётся, но ссылка работать не будет. Наконец, добавьте плагин из листинга~\ref{lst:10-4} и выполните \texttt{mvn clean package}: в \texttt{target} появятся два архива~--- обычный и с суффиксом \texttt{-javadoc}.

\subsection*{Задание 2. Первые тесты, жизненный цикл и исключения}

Добавьте зависимость и плагин из листинга~\ref{lst:10-6} и выполните \texttt{mvn dependency:tree}: убедитесь, что артефакт-агрегат привёл за собой три~--- \texttt{api}, \texttt{params} и \texttt{engine}. Создайте \texttt{FineCalculatorTest} с двумя тестами по схеме AAA (листинг~\ref{lst:10-5}): за 10~дней начисляется 50~рублей, за 1000~дней результат равен максимуму. Запустите все тесты, затем один класс, затем один метод. Сломайте тест намеренно, прочитайте сообщение, загляните в отчёт в каталоге \texttt{target/surefire-reports} и верните значение. Отдельно напишите временный тест \texttt{assertEquals(0.3, 0.1 + 0.2)} без дельты и объясните результат.

Добавьте в основной код класс \texttt{BookShelf} и напишите к нему тест по образцу листинга~\ref{lst:10-7}, дополнив \texttt{setUp} и \texttt{tearDown} печатью в консоль: прочитайте порядок сообщений и объясните письменно, почему поле \texttt{shelf} в разных тестах~--- разные объекты. Затем создайте \texttt{FineCalculatorExceptionTest} по листингу~\ref{lst:10-9} и допишите три теста: с \texttt{assertThrowsExactly} на отказ \texttt{calculateForBooks} при нуле книг, с \texttt{assertDoesNotThrow} на нулевую просрочку и с проверкой границы, где линейный рост упирается в максимум. Уберите из \texttt{calculate} проверку на отрицательное число, прогоните тесты, прочитайте, что скажет JUnit, и верните проверку.

\subsection*{Задание 3. Наборы данных}

Создайте класс \texttt{FineCalculatorParamTest} и перенесите в него три параметризованных теста из листинга~\ref{lst:10-10}, добавив четвёртый: он принимает отрицательные значения (в том числе \texttt{Integer.MIN\_VALUE}) и проверяет, что каждое приводит к исключению. Запустите \texttt{mvn test} и посчитайте, сколько всего прогонов дали эти методы: методов стало меньше, а строк в отчёте больше, потому что каждый набор данных считается отдельным тестом со своим именем.

Допишите в тот же класс тест с \texttt{@MethodSource} по образцу листинга~\ref{lst:10-11}, но для метода \texttt{calculateForBooks}: источник подаёт дни, количество книг, ожидаемую сумму и описание случая, которое служит именем прогона; включите случаи «без просрочки», «три книги по 50~рублей» и «выше максимума штраф не растёт». Замените одно значение на заведомо неверное и убедитесь, что упал ровно один прогон. Наконец, создайте класс \texttt{CardNumberGeneratorTest} со статическим методом, генерирующим восьмизначный номер читательского билета через \texttt{ThreadLocalRandom} и \texttt{String.format}, и тестом с \texttt{@RepeatedTest(10)}, проверяющим длину и состав номера.

\subsection*{Задание 4. Изоляция зависимостей с Mockito}

Добавьте зависимость \texttt{mockito-junit-jupiter} с областью видимости \texttt{test}. Создайте в пакете \texttt{ru.fa.library.model} класс \texttt{Book} (идентификатор, название, признак доступности; конструктор со всеми тремя, геттеры и сеттер признака), а в пакете \texttt{ru.fa.library.service}~--- два интерфейса и сервис из листинга~\ref{lst:10-12}. Выполните \texttt{mvn clean compile} и объясните, почему проект собирается, хотя ни один из интерфейсов не имеет реализации.

Напишите \texttt{LoanServiceTest} по листингу~\ref{lst:10-13}, добавив к приведённому тесту ещё три: занятую книгу выдать нельзя~--- ожидается \texttt{IllegalStateException}, а \texttt{never()} и \texttt{verifyNoInteractions} подтверждают, что ни сохранения, ни письма не было; отсутствующая книга~--- репозиторий возвращает пустой \texttt{Optional}; перехват текста уведомления через \texttt{ArgumentCaptor}. Затем три эксперимента, после каждого возвращая код. Уберите строку настройки \texttt{when} из первого теста и прочитайте ошибку. Замените \texttt{eq($\ldots$)} вокруг адреса на «сырое» значение и получите \texttt{InvalidUseOfMatchersException}. Добавьте в третий тест настройку \texttt{when(repository.save(any())).thenReturn(null)}~--- сервис туда не доходит, заглушка остаётся невостребованной, и строгий режим расширения сообщит об этом исключением \texttt{UnnecessaryStubbingException}; объясните письменно, почему в первом тесте та же строка ошибки не вызвала бы.

\subsection*{Задание 5. Тесты Spring-приложения}

Эта часть выполняется в отдельном проекте на Spring Boot~3.5 с единственной зависимостью Spring Web; тестовый стартер добавляется генератором проекта сам. Координаты: группа \texttt{ru.fa}, артефакт \texttt{student-web}, пакет \texttt{ru.fa.student}, Java~21. Создайте в этом пакете четыре файла: запись \texttt{Student} с компонентами \texttt{id}, \texttt{name}, \texttt{surname}; интерфейс \texttt{StudentService} с методами \texttt{findAll}, \texttt{findById} и \texttt{save}, причём возврат \texttt{null} из \texttt{findById} для отсутствующего студента обязан быть записан в Javadoc, потому что по сигнатуре не угадывается; класс \texttt{StudentServiceImpl} с аннотацией \texttt{@Service}, хранящий студентов в списке в памяти; контроллер \texttt{StudentRestController} с адресом \texttt{/api/students}, отдающий список, одного студента по идентификатору (с кодом~404, если его нет) и создающий нового с кодом~201.

Напишите тест среза \texttt{StudentRestControllerTest} по листингу~\ref{lst:10-14}, добавив ещё два теста: запрос несуществующего студента должен давать код~404, а POST-запрос с телом в формате JSON~--- код~201 и идентификатор в ответе. Обратите внимание на время прогона: контекст поднимается за доли секунды, потому что ни база, ни сервер не запускаются. Уберите аннотацию \texttt{@MockitoBean}, прочитайте ошибку запуска контекста и верните её. Наконец, добавьте интеграционный тест по листингу~\ref{lst:10-15} и прогоните оба класса: в логе появятся строки о старте сервера и случайном порте, а время заметно вырастет. Проверьте разницу уровней на опыте: временно уберите у \texttt{StudentServiceImpl} аннотацию \texttt{@Service}~--- срез веб-слоя этого не заметит, потому что сервис там всё равно макет, а интеграционный тест упадёт ещё на старте контекста.

\subsection*{Результаты занятия}

По итогам занятия у вас должны быть готовы: проект \texttt{library-tests} с документированным классом \texttt{FineCalculator} и файлом \texttt{package-info.java}; сгенерированная документация~--- каталог \texttt{docs} и оба архива из \texttt{target}; тестовые классы \texttt{FineCalculatorTest}, \texttt{BookShelfTest}, \texttt{FineCalculatorExceptionTest}, \texttt{FineCalculatorParamTest}, \texttt{CardNumberGeneratorTest} и \texttt{LoanServiceTest}; проект \texttt{student-web} с тестом среза и интеграционным тестом; вывод \texttt{mvn test} для обоих проектов, где все тесты зелёные.

\kontrolvoprosy

\begin{enumerate}
\item Чем документирующий комментарий отличается от блочного и где он обязан стоять, чтобы утилита его увидела?
\item Что такое дескриптор Javadoc? Чем блочные дескрипторы отличаются от строчных, а \texttt{@throws} от \texttt{@exception}?
\item Чем \texttt{\{@code\}} отличается от \texttt{\{@literal\}}? Что делает \texttt{\{@value\}} и когда необходим \texttt{\{@inheritDoc\}}?
\item Какими опциями команды \texttt{javadoc} задаются каталог вывода и вывод автора? Какие замечания doclint считает ошибками?
\item Опишите пирамиду тестирования и схему AAA. Почему модульных тестов должно быть больше всего?
\item Расшифруйте \eng{FIRST}. Какое из пяти свойств нарушает тест, зависящий от текущей даты?
\item Из каких трёх подпроектов состоит JUnit~5 и за что отвечает каждый? Что такое \texttt{TestEngine}?
\item Сколько экземпляров тестового класса создаёт JUnit при трёх тестовых методах и зачем он так делает?
\item Чем класс \texttt{Assertions} из JUnit~5 отличается от \texttt{Assert} из JUnit~4? Почему при сравнении \texttt{double} нужна дельта и что даёт \texttt{assertAll}?
\item Как проверить выброс исключения в JUnit~5? Что возвращает \texttt{assertThrows} и чем он лучше атрибута \texttt{expected} из JUnit~4?
\item Перечислите источники данных для параметризованного теста. Почему \texttt{@ValueSource} не передаёт \texttt{null}?
\item Объясните разницу между заглушкой, макетом и шпионом. Какое правило действует при смешивании матчеров с обычными значениями? Зачем нужен \texttt{ArgumentCaptor}?
\item Чем \texttt{@WebMvcTest} отличается от \texttt{@SpringBootTest}? Зачем в срезе веб-слоя нужен \texttt{@MockitoBean} и что такое \texttt{MockMvc}?
\item Что показывает покрытие кода и чего оно не показывает? Почему погоня за стопроцентным покрытием вредна?
\end{enumerate}

\testzadaniya

\begin{enumerate}

\vopros{Класс описан фразой «Вычисляет штрафы. Ставка~--- 5~руб. в день.» Что попадёт в сводную таблицу как краткое описание?}
  {«Вычисляет штрафы.»}
  {вся фраза целиком}
  {«Вычисляет штрафы. Ставка~--- 5~руб.»}
  {ничего: для классов краткое описание не формируется}

\vopros{Почему комментарий «увеличиваем i на единицу» над строкой \texttt{i = i + 1} считается вредным?}
  {он замедляет компиляцию}
  {однострочные комментарии не рекомендуются}
  {он мешает собрать документацию метода}
  {он дублирует код, со временем расходится с ним, а человек верит комментарию}

\vopros{Чем дескриптор \texttt{\{@code\}} отличается от \texttt{\{@literal\}}?}
  {\texttt{\{@code\}} подставляет значение константы}
  {\texttt{\{@code\}} создаёт ссылку на класс}
  {оба защищают текст от разбора как HTML, но \texttt{\{@code\}} даёт моноширинный шрифт}
  {\texttt{\{@code\}} работает только в блоке тегов}

\vopros{Какое замечание doclint останавливает генерацию документации?}
  {отсутствующий \texttt{@param} у одного из параметров}
  {отсутствующий \texttt{@return} у метода}
  {битая ссылка в \texttt{\{@link\}}: \texttt{error: reference not found}}
  {отсутствующий тег \texttt{@author} у класса}

\vopros{Какие тесты находятся в основании пирамиды тестирования и почему?}
  {сквозные: они проверяют систему так, как её видит пользователь}
  {интеграционные: они дают баланс скорости и охвата}
  {модульные: они самые быстрые и точнее всех указывают место поломки}
  {ручные: с них начинают проверку любой системы}

\vopros{За что отвечает JUnit Platform?}
  {за выполнение тестов, помеченных аннотацией \texttt{@Test}}
  {за параметризованные тесты и источники данных}
  {за создание заглушек и проверку вызовов}
  {за запуск тестов и связь со средой и сборкой через интерфейс \texttt{TestEngine}}

\vopros{Сколько экземпляров тестового класса создаст JUnit~5 по умолчанию для класса с четырьмя тестами?}
  {один~--- все тесты работают с общим состоянием}
  {два: отдельный для \texttt{@BeforeAll} и общий для тестов}
  {четыре~--- по новому экземпляру на каждый тестовый метод}
  {ни одного: тестовые методы вызываются статически}

\vopros{В каком порядке передаются значения в \texttt{assertEquals}?}
  {сначала фактическое, затем ожидаемое; иначе код не скомпилируется}
  {порядок не важен: сравнение симметрично}
  {сначала фактическое, затем ожидаемое; при перестановке тест всегда падает}
  {сначала ожидаемое, затем фактическое; при перестановке тест работает, но сообщение об ошибке вводит в заблуждение}

\vopros{Что произойдёт при выполнении теста \texttt{double sum = 0.1 + 0.2; assertEquals(0.3, sum);}?}
  {тест пройдёт: JUnit округляет значения типа \texttt{double}}
  {код не скомпилируется: третий аргумент обязателен}
  {тест упадёт, потому что сумма равна \texttt{0.30000000000000004}}
  {тест будет пропущен как недостоверный}

\vopros{Что возвращает метод \texttt{assertThrows} и зачем это нужно?}
  {\texttt{boolean}~--- было ли брошено исключение}
  {строку с сообщением исключения}
  {\texttt{void}; проверить сообщение исключения нельзя}
  {само пойманное исключение~--- его можно исследовать дальше: сообщение, код, причину}

\vopros{Почему массив в \texttt{@ValueSource(strings = $\ldots$)} не может содержать \texttt{null}?}
  {\texttt{null} допустим, но такой прогон молча пропускается}
  {аннотации Java не принимают \texttt{null} в массивах значений; для него есть \texttt{@NullSource}}
  {источник данных работает только с числами}
  {JUnit запрещает прогоны с пустыми аргументами}

\vopros{В чём практическая разница между заглушкой и макетом?}
  {заглушка создаётся вручную, а макет~--- только библиотекой}
  {заглушка применяется к интерфейсам, а макет~--- к классам}
  {заглушка отвечает заранее заданными значениями, а макет ещё и позволяет проверить сами вызовы}
  {заглушка оборачивает настоящий объект, а макет создаётся с нуля}

\vopros{Вызов \texttt{verify(sender).send(адрес, anyString())} с точным значением в первом аргументе даёт исключение. Почему?}
  {метод с типом \texttt{void} нельзя проверять через \texttt{verify}}
  {если хотя бы один аргумент задан матчером, остальные тоже должны быть матчерами: нужен \texttt{eq($\ldots$)}}
  {матчер \texttt{anyString()} не принимает \texttt{null}}
  {порядок аргументов в \texttt{verify} обратный}

\vopros{Что поднимает \texttt{@WebMvcTest} и зачем сервису нужен \texttt{@MockitoBean}?}
  {весь контекст приложения; аннотация лишь ускоряет его запуск}
  {только слой данных; сервис подменяют, чтобы не трогать базу}
  {только веб-слой; сервисов в контексте нет, и бин приходится подменить макетом}
  {веб-слой вместе с сервисами; аннотация только помечает поле}

\vopros{Тест вызывает \texttt{calculate(10)} без единого утверждения, а инструмент показывает стопроцентное покрытие метода. О чём это говорит?}
  {метод протестирован полностью, других тестов не нужно}
  {покрытие фиксирует лишь факт выполнения строк и ничего не говорит о проверке результата}
  {инструмент настроен неверно: покрытие должно быть нулевым}
  {достигнуто покрытие ветвлений, а покрытие строк осталось нулевым}

\end{enumerate}
